\uuid{pbU7}
\titre{Théorème de Bézout}
\auteur{bodin ; croizat ; sacre}
\organisation{exo7}
\licence{CC BY-NC-SA 4.0}
\datecreate{2022-03-08}
\niveau{L1}
\module{Algèbre}
\chapitre{Arithmétique dans Z}
\sousChapitre{Nombres premiers, nombres premiers entre eux}
\difficulte{1}
\theme{Théorème de Bézout}
\exo7qcmid{hilisit.arithmetique.015}

\contenu{
\begin{question}
\qclassification{103.04}
    Soient deux entiers $a,b$ tels que $\pgcd(a,b)=1$.
    Quelles sont les affirmations vraies ?
    \begin{answers}
        \bad{$a$ et $b$ sont des nombres premiers.}
        \good{$a$ et $b$ sont des nombres premiers entre eux.}
        \good{Il existe $u,v\in\Zz$ tels que $au+bv=1$.}
        \good{Il existe $u,v\in\Zz$ tels que $au+bv=2$.}
    \end{answers}
    \begin{explanations}
      $\pgcd(a,b)=1$ est la définition de   $a$ et $b$ sont des nombres premiers entre eux.
      Le théorème de Bézout affirme qu'il existe $u,v\in\Zz$ tels que $au+bv=1$.
      En multipliant cette égalité par $2$, on obtient $a(2u)+b(2v)=2$.
    \end{explanations}
\end{question}
}
